\chapter{Produits raffinés et cracks}\label{ch:m3:refined-products-and-cracks}

Une raffinerie de la côte américaine du Golfe achète du brut et vend de l'essence et du gazole. Elle
ne se soucie guère que le pétrole coûte soixante ou cent dollars~; ce qui compte pour elle, c'est
l'écart entre ce que rapportent ses produits et ce que lui coûte son brut. Cet
écart porte un nom sur tous les écrans de l'énergie, le crack, d'après le craquage
des molécules lourdes en molécules légères. En août 2026, la version la plus courante de cet
écart s'établissait en moyenne à environ 65~\$ le baril de brut à New York Harbor, soit trois fois sa
moyenne des vingt années précédentes, et un raffineur pouvait verrouiller l'essentiel de la marge
d'un trimestre en un après-midi, avec trois contrats à terme. Ce chapitre traite des
produits, des spreads qui valorisent l'activité de raffinage, de l'arbitrage qui déplace
les produits d'une région à l'autre, ainsi que des saisons et des réglementations qui façonnent l'ensemble.

\section{La raffinerie et sa gamme de produits raffinés}

La distillation sépare le brut selon son point d'ébullition~: les gaz et le naphta en tête,
le kérosène et le gazole au milieu, le fioul lourd et le résidu en bas.
Les unités de conversion (craqueurs catalytiques, hydrocraqueurs, cokéfaction) découpent les
fractions lourdes en fractions plus légères et plus valorisées~; les unités de traitement éliminent le
soufre~; l'assemblage combine les composants en produits qui satisfont une spécification.

\begin{definition}[Gamme de produits raffinés]\label{def:m3:refined-products-and-cracks:slate}
La \emph{gamme de produits raffinés}\index{gamme de produits raffinés} d'une raffinerie est l'ensemble des produits qu'elle
fabrique et leur part dans sa production, par baril de brut traité.
\end{definition}

Une raffinerie américaine tire environ 19 à 20~gallons d'essence et 11 à 13~gallons
de distillat à très faible teneur en soufre tiré d'un baril de brut de 42~gallons~; la gamme varie
avec le brut, la saison et les unités dont dispose la raffinerie. La production totale
dépasse en général 42~gallons, car les produits sont moins denses que le brut~:
les raffineurs parlent du gain de traitement.

\begin{omfigure}
\centering
\begin{tikzpicture}[font=\small,
  box/.style={draw,thick,rounded corners=2pt,align=center,minimum height=0.8cm},
  arr/.style={-{Stealth[length=2.2mm]},thick}]
\node[box,draw=omDef,fill=omDef!8,minimum width=2.2cm] (crude) at (0,0) {pétrole brut\\42 gallons};
\node[box,draw=omThm,fill=omThm!8,minimum width=1.8cm,minimum height=3.6cm] (col) at (3.6,0) {distillation\\et\\conversion};
\draw[arr] (crude) -- (col);
\node[box,draw=omProp,fill=omProp!8,minimum width=4.4cm,anchor=west] (g) at (6.0,1.35) {essence~: environ 19--20 gallons};
\node[box,draw=omProp,fill=omProp!8,minimum width=4.4cm,anchor=west] (d) at (6.0,0.2) {distillat~: environ 11--13 gallons};
\node[box,draw=omExo,fill=omExo!8,minimum width=4.4cm,anchor=west] (o) at (6.0,-0.95) {carburéacteur, fioul, autres};
\draw[arr] (col.east |- g) -- (g);
\draw[arr] (col.east |- d) -- (d);
\draw[arr] (col.east |- o) -- (o);
\node[font=\footnotesize,anchor=north] at (3.6,-2.0) {les produits réunis dépassent 42 gallons~: gain de traitement};
\end{tikzpicture}

\omcaption{Un baril de brut passant par une raffinerie américaine. Les parts varient selon le brut, la saison et
la configuration de la raffinerie. Source des volumes d'essence et de distillat~: US Energy
Information Administration. Schéma.}\label{fig:m3:refined-products-and-cracks:slate}
\end{omfigure}

\section{Essence, distillats et fioul}

\begin{definition}[Distillats moyens]\label{def:m3:refined-products-and-cracks:middle}
Les \emph{distillats moyens}\index{distillats moyens} sont les produits issus de la
partie médiane de la plage de distillation~: le kérosène et le carburéacteur, le gazole et le fioul domestique, ainsi que les
gazoles dont ils sont tirés.
\end{definition}

Chaque produit a ses propres indices de référence. Aux États-Unis, les contrats à terme sont
le RBOB (gasoline blendstock) et le gazole à très faible teneur en soufre (ULSD), tous deux livrés à
New York Harbor par lots de 42\,000~gallons (1\,000~barils) et cotés en dollars
par gallon~; en Europe, le gazole à faible teneur en soufre, livré dans la zone Amsterdam--Rotterdam--Anvers
par lots de 100~tonnes et coté en dollars par tonne. Une tonne de gazole de
masse volumique \qty{0.845}{kg/l} représente 1.1835~mètres cubes, soit 7.44~barils~: pour la comparer au brut,
il suffit de diviser le prix à la tonne par 7.44.

\begin{definition}[Plafond de soufre des carburants marins]\label{def:m3:refined-products-and-cracks:imo}
Le \emph{plafond de soufre des carburants marins}\index{plafond de soufre des carburants marins} est la
limite fixée par l'Organisation maritime internationale sur la teneur en soufre du carburant utilisé par
les navires~: 0.50~\% en masse depuis le 1\up{er} janvier 2020 en dehors des zones de contrôle des émissions (contre
3.50~\% auparavant), et 0.10~\% à l'intérieur de ces zones depuis 2015.
\end{definition}

Ce plafond a déplacé la demande du fioul à forte teneur en soufre, ce résidu lourd que les navires avaient
brûlé pendant des décennies, vers des mélanges à faible teneur en soufre et vers le gazole marin. C'est
l'exemple le plus net d'une réglementation qui a réévalué un produit du jour au lendemain~: la valeur
du résidu d'une raffinerie dépendait d'un jour à l'autre de la présence ou non d'unités
capables de le désulfurer.

\section{Crack spreads et marges de raffinage}

\begin{definition}[Crack spread, crack spread 3-2-1]\label{def:m3:refined-products-and-cracks:crack}
Un \emph{crack spread}\index{crack spread} est l'écart entre la valeur des
produits raffinés et le coût du brut dont ils sont issus, selon un rapport de volumes fixe
par baril de brut. Le \emph{crack spread 3-2-1}\index{crack spread 3-2-1}
vend deux barils d'essence et un de gazole pour trois barils de brut~:
\[
\mathrm{C}_{321} = \frac{2\,P_{\mathrm{gas}} + P_{\mathrm{dist}} - 3\,P_{\mathrm{crude}}}{3},
\]
tous les prix en dollars par baril (les prix des produits, exprimés par gallon, multipliés par 42).
\end{definition}

\begin{example}[Un crack à partir de trois cotations]\label{ex:m3:refined-products-and-cracks:crack}
Avec un brut à \qty{83.90}{\$/\barrel}, l'essence à \qty{3.2132}{\$/gal} et le gazole à
\qty{4.2530}{\$/gal} (moyennes d'août 2026 des données utilisées ci-dessous), les produits
valent \qty{134.95}{\$/\barrel} et \qty{178.63}{\$/\barrel}, et le crack 3-2-1 vaut
\qty{65.61}{\$/\barrel}. Séparément, le crack de l'essence vaut \qty{51.06}{} et le crack du gazole
vaut \qty{94.73}{\$/\barrel}.
\end{example}

Un crack est une raffinerie stylisée~: une raffinerie réelle a sa propre gamme et ses propres coûts.

\begin{definition}[Marge brute de raffinage]\label{def:m3:refined-products-and-cracks:grm}
La \emph{marge brute de raffinage}\index{marge brute de raffinage} d'une raffinerie est la
valeur des produits qu'elle obtient d'un baril d'un brut donné, aux prix de marché
et aux rendements de ces produits (gain de traitement compris), diminuée du prix de ce brut. La marge nette
déduit en outre les coûts d'exploitation.
\end{definition}

\begin{proposition}[Un crack est un portefeuille de contrats à terme]\label{prop:m3:refined-products-and-cracks:hedge}
Un raffineur qui traite $Q$ barils de brut sur une période et vend ses produits dans le
rapport $2:1$ fixe le \omterm{def:m3:refined-products-and-cracks:crack}{crack spread 3-2-1} en vigueur aujourd'hui en achetant $Q$ barils de brut
à terme et en vendant $2Q/3$ de contrats à terme sur l'essence et $Q/3$ sur le gazole pour cette période.
Sa marge ne varie alors qu'en fonction de l'écart entre ses propres rendements, ses localisations et
ses qualités et ceux des contrats.
\end{proposition}

\begin{proof}
Sur le volume couvert, la marge physique du raffineur par baril vaut $\frac{1}{3}(2P_g + P_d
- 3P_c)$ aux prix de la période~; le gain sur les contrats à terme vaut $\frac{1}{3}(2(F_g - P_g) + (F_d -
P_d) - 3(F_c - P_c))$ par baril si les prix physiques égalent les cours à terme à
la livraison. La somme est le crack aux cours à terme du jour.
\end{proof}

\begin{omfigure}
\begin{tikzpicture}
\begin{axis}[width=11cm,height=5cm,
  ylabel={\$ par baril de brut},
  tick label style={font=\small},label style={font=\small},
  xmin=2006.5,xmax=2026.75,ymin=0,ymax=100,grid=major,grid style={gray!25},
  xtick={2008,2010,2012,2014,2016,2018,2020,2022,2024,2026},xticklabel style={/pgf/number format/1000 sep={}},
  legend columns=3,legend style={font=\footnotesize,at={(0.5,-0.18)},anchor=north,draw=none,fill=none,/tikz/every even column/.append style={column sep=8pt}},
  table/col sep=comma]
\addplot[omDef,very thick] table[x=t,y=crack]{figdata/markets-3/03-refined-products-and-cracks/crack321.csv};
\addplot[omProp,thin] table[x=t,y=gas]{figdata/markets-3/03-refined-products-and-cracks/crack321.csv};
\addplot[omThm,thin,dashed] table[x=t,y=diesel]{figdata/markets-3/03-refined-products-and-cracks/crack321.csv};
\legend{crack spread 3-2-1,crack de l'essence,crack du gazole}
\end{axis}
\end{tikzpicture}

\omcaption{Cracks de New York Harbor face au WTI, moyennes mensuelles des prix spot
quotidiens, de juillet 2006 à août 2026. Les sommets sont ceux de mai 2022 (\$62.68) et d'août 2026
(\$65.61, le plus haut)~; le creux est celui de novembre 2009 (\$5.51). Données~: séries FRED
DCOILWTICO, DGASNYH et DDFUELNYH (US Energy Information
Administration).}\label{fig:m3:refined-products-and-cracks:crack}
\end{omfigure}

Le crack s'est établi en moyenne à \qty{21.20}{\$/\barrel} sur la période
(\cref{fig:m3:refined-products-and-cracks:crack}). Il a progressé en 2022, après que l'invasion à grande échelle de l'Ukraine par la Russie a mis en péril la source
de la moitié des importations européennes de gazole (l'Union européenne a interdit le gazole russe transporté par mer
au début de 2023), et de nouveau en 2026, lorsque la perturbation du Golfe décrite au \cref{ch:m3:crude-oil} a retiré les produits raffinés aussi bien que le brut
du marché~; dans les deux épisodes, le crack du gazole a dépassé de loin celui de l'essence.

\begin{dated}{2026-09}{Capacité de raffinage aux États-Unis}\label{dat:m3:refined-products-and-cracks:capacity}
La capacité de distillation atmosphérique en service aux États-Unis s'élevait à 18.2 millions de barils par jour calendaire
au 1\up{er} janvier 2026, en baisse d'environ 250\,000 barils par jour (environ 1~\%) sur un an,
pour 130 raffineries en service (301 en 1982).
\end{dated}

\section{Arbitrage régional}

Un même produit se négocie à des prix différents selon les lieux, et circule
entre eux par pipeline ou par navire lorsque l'écart couvre le coût du déplacement.

\begin{definition}[Fenêtre d'arbitrage]\label{def:m3:refined-products-and-cracks:arb}
Une \emph{fenêtre d'arbitrage}\index{fenêtre d'arbitrage} pour un produit entre deux régions
est ouverte lorsque son prix à destination dépasse son prix à l'origine de plus
que le coût de son transport (fret, assurance, pertes, droits et financement pendant le
voyage), et fermée dans le cas contraire.
\end{definition}

La côte américaine du Golfe raffine bien plus qu'elle ne consomme et expédie le surplus vers la
côte Est par pipeline et par pétrolier, ainsi qu'à l'étranger. L'essence de New York Harbor se négocie donc
habituellement au-dessus de l'essence de la côte du Golfe~: de 6.0 cents par gallon en moyenne dans les
données mensuelles depuis 2006, et en dessous dans seulement 35 mois sur 242
(\cref{fig:m3:refined-products-and-cracks:arb}). Lorsque l'écart passe sous le
coût de transport d'un gallon, le flux ralentit~; lorsqu'il dépasse ce coût, davantage de produit circule.
Les \omterm{def:m3:refined-products-and-cracks:arb}{fenêtres d'arbitrage} relient les marchés régionaux de produits, et ce sont elles qui
les déconnectent lorsqu'un pipeline se rompt ou qu'un canal ferme.

\begin{omfigure}
\begin{tikzpicture}
\begin{axis}[width=11cm,height=4.4cm,
  ylabel={cents par gallon},
  tick label style={font=\small},label style={font=\small},
  xmin=2006.5,xmax=2026.75,ymin=-40,ymax=70,grid=major,grid style={gray!25},
  xtick={2008,2010,2012,2014,2016,2018,2020,2022,2024,2026},xticklabel style={/pgf/number format/1000 sep={}},
  table/col sep=comma]
\addplot[omDef,thick] table[x=t,y=cents]{figdata/markets-3/03-refined-products-and-cracks/arb.csv};
\draw[gray] (axis cs:2006.5,0) -- (axis cs:2026.75,0);
\end{axis}
\end{tikzpicture}

\omcaption{New York Harbor moins côte américaine du Golfe, pour l'essence conventionnelle, moyennes mensuelles
des prix spot quotidiens, de juillet 2006 à août 2026. Le spread valorise le déplacement du
centre de raffinage vers la côte de consommation. Données~: séries FRED DGASNYH et DGASUSGULF (US
Energy Information Administration).}\label{fig:m3:refined-products-and-cracks:arb}
\end{omfigure}

\section{Saisonnalité et spécifications}

Les produits ont leurs saisons. La demande d'essence culmine pendant la saison estivale de conduite dans l'hémisphère Nord,
la demande en \omterm{def:m3:refined-products-and-cracks:middle}{distillats moyens} pendant le chauffage hivernal~; les raffineries programment l'entretien au printemps et
à l'automne, entre les pics. Les spécifications changent elles aussi avec la saison.

\begin{definition}[Pression de vapeur Reid]\label{def:m3:refined-products-and-cracks:rvp}
\emph{pression de vapeur Reid}\index{pression de vapeur Reid} (RVP) est la mesure standard
de la volatilité de l'essence, en livres par pouce carré~; un RVP plus bas signifie une évaporation moindre
et moins d'émissions de composés organiques volatils.
\end{definition}

L'essence d'été américaine, vendue du 1\up{er} juin au 15 septembre, ne peut dépasser un RVP de
\qty{9.0}{psi}, et \qty{7.8}{psi} dans certaines zones. La qualité d'été exige moins de composants légers bon marché
comme le butane et coûte plus cher à produire, si bien que le contrat à terme sur essence
change de qualité au printemps et son cours fait un bond~: le règlement de marché du NYMEX plafonne le RVP du RBOB livré à \qty{13.5}{psi} en mars et
\qty{7.4}{psi} d'avril à la mi-septembre.
\Cref{fig:m3:refined-products-and-cracks:season} montre l'effet net sur les cracks~:
De 2010 à 2025, le crack de l'essence s'est situé en moyenne à \qty{2.83}{\$/\barrel} au-dessus de sa moyenne en
mai et à \qty{3.83}{\$/\barrel} en dessous de sa moyenne en janvier~; le crack du gazole a culminé en
novembre, \qty{4.15}{\$/\barrel} au-dessus de sa moyenne.

\begin{omfigure}
\begin{tikzpicture}
\begin{axis}[width=11cm,height=4.8cm,ybar,bar width=5pt,
  xlabel={mois calendaire},ylabel={\$/bbl par rapport à la moyenne},
  tick label style={font=\small},label style={font=\small},
  xmin=0.4,xmax=12.6,ymin=-5,ymax=5,grid=major,grid style={gray!25},
  xtick={1,2,3,4,5,6,7,8,9,10,11,12},
  legend columns=2,legend style={font=\footnotesize,at={(0.5,-0.42)},anchor=north,draw=none,fill=none,/tikz/every even column/.append style={column sep=8pt}},
  table/col sep=comma]
\addplot[fill=omProp!60,draw=omProp] table[x=month,y=gas]{figdata/markets-3/03-refined-products-and-cracks/season.csv};
\addplot[fill=omThm!60,draw=omThm] table[x=month,y=diesel]{figdata/markets-3/03-refined-products-and-cracks/season.csv};
\legend{crack de l'essence,crack du gazole}
\end{axis}
\end{tikzpicture}

\omcaption{Saisonnalité des cracks de New York Harbor~: la moyenne par mois calendaire moins la
moyenne globale, de 2010 à 2025. L'essence est chère pendant la saison de conduite, le gazole
pendant la saison de chauffe. Données~: voir la \cref{fig:m3:refined-products-and-cracks:crack}.}\label{fig:m3:refined-products-and-cracks:season}
\end{omfigure}

\section{Tutoriel~: cracks à partir des prix, et la couverture}\label{tut:m3:refined-products-and-cracks:crack}

\begin{tutorial}
\textbf{Objectif.} Calculer les cracks à partir des prix des produits et du brut dans leurs propres unités, et
dimensionner, par contrats à terme, la couverture de la \cref{prop:m3:refined-products-and-cracks:hedge}.
\textbf{État final~:} \cref{fig:m3:refined-products-and-cracks:crack,fig:m3:refined-products-and-cracks:season}
et les lots du problème du week-end.
\begin{tutsteps}
\item \textbf{Le crack.} Tout rapport, avec la vérification que les produits et le brut s'équilibrent.
\omcode{code/firm/cracks/firm_cracks.py}{27}{38}{Un crack spread et le crack 3-2-1.}\label{lst:m3:refined-products-and-cracks:crack}
\item \textbf{La couverture.} Brut acheté, produits vendus, par lots entiers.
\omcode{code/firm/cracks/firm_cracks.py}{47}{54}{Lots de contrats à terme qui figent un crack sur un débit.}\label{lst:m3:refined-products-and-cracks:lots}
\item \textbf{Exécuter} \texttt{m3\_cracks.crack\_series()}, \texttt{seasonality()},
\texttt{refinery\_hedge()} et \texttt{fig\_cracks.py}.
\end{tutsteps}
\textbf{Pour aller plus loin.} Remplacer le 3-2-1 par un 5-3-2 et par un crack européen sur
le Brent et le gazole, en tonnes~; décaler les produits d'un mois (le temps que le brut met à
devenir produit) et observer comment la volatilité du crack change.
\end{tutorial}

\section{Construction~: le calculateur de crack}\label{bld:m3:refined-products-and-cracks:cracks}

\begin{build}
\bfield{Rôle} La mini-entreprise négocie des marges de raffinage pour les couvertures des raffineurs et
pour son propre portefeuille~: elle a besoin de cracks dans tout rapport et toute unité, de marges de raffinage à partir des rendements,
de tailles de couverture en lots, et d'un test indiquant si une \omterm{def:m3:refined-products-and-cracks:arb}{fenêtre d'arbitrage} est ouverte.
\bfield{Interface} \texttt{per\_gallon\_to\_per\_barrel}~; \texttt{barrels\_per\_tonne(density)}~;
\texttt{per\_tonne\_to\_per\_barrel}~; \texttt{crack(crude, products, ratio, crude\_barrels)}~;
\texttt{three\_two\_one}~; \texttt{gross\_refining\_margin(crude, yields, products)}~;
\texttt{hedge\_lots(throughput\_bpd, days, ratio, crude\_barrels)}~;
\texttt{arbitrage\_open}.
\bfield{Règles} Tout est converti en dollars par baril avant d'être soustrait~; un
crack dont les barils de produits ne sont pas égaux à ses barils de brut est rejeté~; les lots sont
arrondis à des contrats entiers et les lots de la couverture somment à zéro.
\bfield{Tests d'acceptation} \texttt{code/firm/cracks/tests/}~: 7.44 barils par tonne de
gazole~; le 3-2-1 comme crack~; une marge avec gain de traitement~; une couverture équilibrée.
\bfield{Pour aller plus loin} Un crack sur des contrats à terme de mois différents (brut ce mois-ci, produits
le mois suivant)~; un programme linéaire de raffinerie qui choisit la \omterm{def:m3:refined-products-and-cracks:slate}{gamme de produits raffinés} à partir des prix des produits.
\end{build}

\begin{omsources}
\item US Energy Information Administration~: \emph{Frequently Asked Questions} (products
from a barrel)~; \emph{Refinery Capacity Report} et \emph{Today in Energy}, ``U.S.
refining capacity decreased during 2025'' (2026)~; série des prix spot via FRED.
\item International Maritime Organization, ``IMO 2020 --- cutting sulphur oxide emissions''.
\item US Environmental Protection Agency, ``Gasoline Reid vapor pressure''.
\item NYMEX Rulebook, chapitres 191 (RBOB gasoline) et 150 (NY Harbor ULSD) (copies de l'Internet Archive)~; spécification du contrat à terme
ICE Low Sulphur Gasoil.
\end{omsources}

\section{Exercices}

\begin{exercise}[$\star$]\label{exo:m3:refined-products-and-cracks:1}
Le brut vaut \qty{75}{\$/\barrel}, l'essence \qty{2.40}{\$/gal}, le gazole \qty{2.80}{\$/gal}. Calculer
le crack 3-2-1.
\end{exercise}

\begin{exercise}[$\star$]\label{exo:m3:refined-products-and-cracks:2}
Le gazole se négocie à \qty{900}{\$/t}. Quel est ce prix par baril, à la densité du
contrat ICE~?
\end{exercise}

\begin{exercise}[$\star$]\label{exo:m3:refined-products-and-cracks:3}
Pourquoi l'essence se négocie-t-elle plus haut en mai qu'en janvier, et pourquoi son cours à terme
grimpe-t-il au printemps~?
\end{exercise}

\begin{exercise}[$\star\star$]\label{exo:m3:refined-products-and-cracks:4}
Les rendements d'une raffinerie sont de 47~\% pour l'essence, 30~\% pour le gazole et 30~\% pour d'autres produits valant
\qty{60}{\$/\barrel}. Avec un brut à \qty{80}{}, l'essence à \qty{100}{} et le gazole à
\qty{110}{\$/\barrel}, quelle est sa marge brute~? Pourquoi les rendements totalisent-ils plus de 100~\%~?
\end{exercise}

\begin{exercise}[$\star\star$]\label{exo:m3:refined-products-and-cracks:5}
L'essence de New York vaut \qty{2.62}{\$/gal}, celle du Golfe \qty{2.50}{\$/gal}, et transporter
un gallon coûte \qty{0.08}{\$}. La fenêtre est-elle ouverte~? À quel cours du Golfe se refermerait-elle~?
\end{exercise}

\begin{exercise}[$\star\star$]\label{exo:m3:refined-products-and-cracks:6}
Combien de lots CL, RB et HO figent un crack 3-2-1 sur 60\,000 barils par jour pendant 30 jours~?
\end{exercise}

\begin{exercise}[$\star\star\star$]\label{exo:m3:refined-products-and-cracks:7}
\emph{Codage.} Avec \texttt{crack\_series}, trouver le mois où le crack 3-2-1 est le plus élevé en
2022, ainsi que sa valeur et le crack moyen sur l'ensemble des mois.
\end{exercise}

\begin{exercise}[$\star\star\star$]\label{exo:m3:refined-products-and-cracks:8}
\emph{Repérer le défaut.} « Notre raffinerie est couverte~: nous avons acheté des contrats à terme sur le brut pour
tout notre débit.~»
\end{exercise}

\section{Problème~: couvrir le trimestre d'une raffinerie}

\begin{problem}[{Problème du week-end --- figer une marge record}]\label{pb:m3:refined-products-and-cracks:1}
À la fin d'août 2026, une raffinerie traitant 90\,000 barils par jour veut
figer sa marge pour le trimestre suivant (92 jours). Prendre les moyennes d'août comme
prix à terme~: WTI \qty{83.8976}{\$/\barrel}, essence \qty{3.2132}{\$/gal}, gazole
\qty{4.2530}{\$/gal}. Sa \omterm{def:m3:refined-products-and-cracks:slate}{gamme de produits raffinés} est proche de 2 barils d'essence pour 1 de gazole.

\medskip\noindent\textbf{Partie I --- La marge.}
\begin{enumerate}
\item Donner les prix de l'essence et du gazole par baril.
\item Donner le \omterm{def:m3:refined-products-and-cracks:crack}{crack spread 3-2-1}.
\item Comment se compare-t-il à la moyenne depuis 2006~?
\item Pourquoi les \omterm{def:m3:refined-products-and-cracks:middle}{distillats moyens} ont-ils été beaucoup plus chers que l'essence~?
\item Quel a été le précédent sommet du crack, et quand~?
\end{enumerate}

\medskip\noindent\textbf{Partie II --- La couverture.}
\begin{enumerate}[resume]
\item Combien de barils de brut la raffinerie va-t-elle traiter au cours du trimestre~?
\item Combien de lots CL achète-t-elle~?
\item Combien de lots RB et HO vend-elle~?
\item Quelle marge, en dollars, la couverture fige-t-elle~?
\item Si le crack tombe à \$25 à la livraison, que gagnent ou perdent les contrats à terme~?
\end{enumerate}

\medskip\noindent\textbf{Partie III --- Ce qui n'est pas couvert.}
\begin{enumerate}[resume]
\item Le brut de la raffinerie n'est pas le WTI de Cushing. Quel risque subsiste~?
\item Ses produits sont vendus sur la côte du Golfe, et non à New York. Quel risque subsiste~?
\item Elle produit 3 barils d'essence pour 2 de gazole. Comment la couverture doit-elle changer~?
\item Un arrêt imprévu immobilise la raffinerie pendant deux semaines. Que devient la couverture~?
\item Quelle marge la raffinerie doit-elle déposer sur la couverture si le crack s'élargit encore~?
\end{enumerate}

\medskip\noindent\textbf{Partie IV --- Jugement.}
\begin{enumerate}[resume]
\item Pourquoi la raffinerie ne couvrirait-elle que la moitié du trimestre~?
\item Qui prend l'autre côté de la couverture de crack d'un raffineur~?
\item Pourquoi un crack record peut-il être un bon moment pour couvrir, et pourtant un moment difficile~?
\item Énoncer le \emph{résultat nommé}~: les lots et la marge que la couverture fige.
\item En une phrase~: que vend une raffinerie~?
\end{enumerate}
\end{problem}

\section{Questions d'entretien}

\begin{interviewq}[$\star$ \iqroles{trader}]\label{iq:m3:refined-products-and-cracks:1}
Qu'est-ce qu'un \omterm{def:m3:refined-products-and-cracks:crack}{crack spread 3-2-1}, et pourquoi ces nombres~?
\end{interviewq}

\begin{interviewq}[$\star$ \iqroles{trader, chercheur}]\label{iq:m3:refined-products-and-cracks:2}
Le gazole est coté par tonne et le brut par baril. Comment les comparez-vous~?
\end{interviewq}

\begin{interviewq}[$\star\star$ \iqroles{chercheur}]\label{iq:m3:refined-products-and-cracks:3}
Comment testeriez-vous si le crack de l'essence est saisonnier, et le négocieriez-vous si c'est le cas~?
\end{interviewq}

\begin{interviewq}[$\star\star$ \iqroles{trader}]\label{iq:m3:refined-products-and-cracks:4}
Les cracks du gazole triplent en un mois alors que ceux de l'essence bougent à peine. Qu'est-ce qui pourrait en être la cause~?
\end{interviewq}

\begin{interviewq}[$\star\star$ \iqroles{risque}]\label{iq:m3:refined-products-and-cracks:5}
Un raffineur couvre son crack un an à l'avance. Lister les risques que la couverture ne couvre pas.
\end{interviewq}

\begin{interviewq}[$\star\star\star$ \iqroles{développeur, chercheur}]\label{iq:m3:refined-products-and-cracks:6}
Concevoir un rapport quotidien des marges de raffinage pour vingt raffineries, chacune
avec son propre brut et sa propre gamme de produits, à partir des prix de marché.
\end{interviewq}
